\chapter{INSTRUMENTO PROPOSTO DE AVALIAÇÃO}
\label{chapter-E}
\section{Protocolo ainda não executado}
\label{section-E-1}
Este instrumento é uma proposta para avaliação futura e não constitui resultado da pesquisa. Antes de aplicá-lo, definir participantes, consentimento, tratamento de dados e exigências institucionais cabíveis. Utilizar somente casos fictícios e informar a natureza demonstrativa do produto. Não coletar senhas, links de autenticação, número real de processo ou nome de parte.

As tarefas propostas são: calcular o caso civil da Capital; explicar a diferença entre publicação e primeiro dia computado; identificar o aviso de calendário parcial; comparar Santos com Capital; salvar e localizar uma memória; agendar e cancelar um lembrete. Observar conclusão, pedidos de ajuda, interpretação de erros e compreensão do significado de aceite do provedor.

O formulário pode registrar perfil de experiência sem identificação nominal, tarefa, conclusão sem ajuda/com ajuda/não concluída, dificuldade relatada e sugestão. Perguntar se a pessoa compreendeu os limites e onde procuraria os fundamentos. A revisão jurídica deve ser conduzida em trilha separada, por profissional identificado, com casos, fonte e parecer registrados; satisfação com a interface não aprova uma regra.

\section{Critérios futuros de análise}
\label{section-E-2}
Apresentar resultados somente após aplicação real, com número de participantes, método, condições e limites. Se forem usadas escalas padronizadas, citar o instrumento e seguir seu procedimento. Não preencher retrospectivamente pontuações SUS, percentuais de sucesso ou comentários fictícios. Falhas observadas devem gerar requisitos e novos testes, preservando a distinção entre evidência e hipótese.
